% !TeX root = ../../main.tex
\subsection{最上級を使った構文}
    \begin{theorembox}{\textbf{最上級の基本構文}}
        \begin{itemize}[nosep, leftmargin=1.5em]
            \item the + 最上級 + \fitblankbf[]{of} + (\fitblankbf[]{数を表す語句})
            \item the + 最上級 + \fitblankbf[]{in} + (\fitblankbf[]{場所・範囲を表す語句})
        \end{itemize}
        \theoremmeaning{「～の中で \fitblankbf[]{一番} (形容詞 [ 副詞 ]) である」}
    \end{theorembox}

    \begin{theorembox}{\textbf{one of the + 最上級}}
        \fitblankbf[]{one of the} + 最上級 + (\fitblankbf[]{名詞の複数形})
        \theoremmeaning{「\fitblankbf[]{最も} (形容詞 [ 副詞 ]) な (名詞) の \fitblankbf[]{1つ}」}
    \end{theorembox}

    \begin{theorembox}{\textbf{the + 序数 + 最上級}}
        \theoremline{the + (\fitblankbf[]{序数}) + 最上級}{「～\fitblankbf[]{番目} に (形容詞 [ 副詞 ]) な … 」}
    \end{theorembox}

    ※序数とは \uwave{first, second, third などの数} のことである。

    ※副詞の最上級には \fitblankbf[]{the} を付けない場合もある。

\subsection{練習問題}
    \begin{enumerate}[leftmargin=*, nosep, itemsep=0.65em]
        \item Mount Everest is the highest mountain in the world.
        \dialogueblank{エベレスト山は世界で一番高い山だ。}
        \item これは3つのうちで一番安い時計だ。
        \dialogueblank{This is the cheapest of the three watches.}
        \item 彼は日本で最も有名な野球選手の1人だ。
        \dialogueblank{He is one of the most famous baseball players in Japan.}
        \item 札幌は日本で5番目に大きな都市だ。
        \dialogueblank{Sapporo is the fifth largest city in Japan.}
        \item 彼はクラスの中で最も上手に英語を話す。
        \dialogueblank{He speaks English best in the class.}
        \item Which planet is the closest to the Sun?
        \dialogueblank{どの惑星が太陽に最も近いですか？}
    \end{enumerate}

\newpage
